% !TeX program = xelatex
% =====================================================================
%  完美编织（Perfect Braid）的定义与评判标准
%  问题四 定义文档（可独立编译，亦可并入主论文）
%  编译：xelatex -> xelatex
% =====================================================================
\documentclass[12pt, a4paper]{ctexart}

\usepackage[top=2.5cm, bottom=2.5cm, left=2.6cm, right=2.6cm]{geometry}
\setCJKmainfont[AutoFakeBold=2.5, AutoFakeSlant=0.2]{SimSun}
\setCJKsansfont[AutoFakeBold=2.5]{SimHei}
\setCJKmonofont{SimHei}
\setmainfont{Times New Roman}
\setsansfont{Arial}

\usepackage{amsmath, amssymb, amsfonts}
\usepackage{amsthm}
\usepackage{bm}
\usepackage{booktabs}
\usepackage{array}
\usepackage{multirow}
\usepackage{makecell}
\usepackage{adjustbox}
\usepackage{enumitem}
\usepackage{textcomp}
\usepackage{tcolorbox}
\usepackage[hidelinks]{hyperref}

\linespread{1.4}
\setlength{\parindent}{2em}
\setlength{\parskip}{2pt plus 1pt}
\emergencystretch=3em
\tolerance=2000

% --- 定理类环境 ---
\newtheorem{definition}{定义}[section]
\newtheorem{proposition}{命题}[section]
\newtheorem{theorem}{定理}[section]
\newtheorem{remark}{注}[section]

\newtcolorbox{keybox}[1][]{
  colback=blue!4, colframe=blue!55!black,
  boxrule=0.7pt, arc=2pt, left=6pt, right=6pt, top=5pt, bottom=5pt,
  fonttitle=\bfseries\sffamily, title={#1}
}

% --- 单位宏（\Omega 不可放进 \mathrm，否则 Times New Roman 渲染为空白方框）---
\newcommand{\U}[1]{\ensuremath{\,\mathrm{#1}}}
\newcommand{\Umm}{\ensuremath{\,\mathrm{mm}}}
\newcommand{\Umohm}{\ensuremath{\,\mathrm{m}\Omega}}
\newcommand{\Uohm}{\ensuremath{\,\mathrm{m}\Omega/\mathrm{m}}}
\newcommand{\UdegC}{\ensuremath{\,\text{\textdegree}\mathrm{C}}}
\newcommand{\Rac}{R_{\mathrm{ac}}}
\newcommand{\Rdc}{R_{\mathrm{dc}}}

\title{\bfseries 完美编织（Perfect Braid）的定义与评判标准\\[6pt]
\large 三维拓扑编织 Litz 线电磁优化设计 —— 问题四理论基础}
\author{}
\date{}

\begin{document}
\maketitle
\thispagestyle{empty}

\begin{keybox}[文档定位]
本文档给出「完美换位 / 完美编织」的严格数学定义、可直接计算的形式化判据，
以及一个可参数化的编织模型（含母轨迹参数方程与生成变换）。
全部定义域、参数与约束条件均显式给出，判据可用统一流程计算，
使不同编织方案之间可以\textbf{客观排序}。
文中的命题均已数值验证（见第 \ref{sec:verify} 节）。
\end{keybox}

\section{目的与范围}

\subsection{要解决的问题}

Litz 线中每股细丝的高频损耗取决于它所经历的磁场历史。邻近效应损耗近似正比于
该丝所处位置的磁场强度平方，而磁场强度沿径向单调变化（越靠外越强）。
因此，\textbf{各丝的径向停留分布是否公平}，决定了电流与损耗能否被均分。

本文档要形式化回答三个问题：

\begin{enumerate}[label=(\roman*), leftmargin=2.4em, itemsep=2pt]
\item \textbf{什么是「完美编织」}——给出无歧义的数学定义；
\item \textbf{如何度量「接近完美」}——给出可计算、可比较的量化判据；
\item \textbf{如何构造与检验}——给出参数化模型与验证流程。
\end{enumerate}

\subsection{不涉及的内容}

本文档只处理\textbf{几何与拓扑层面}的公平性判据，不涉及具体电磁仿真的数值结果；
判据的物理有效性由仿真结果另行验证。

\section{记号、坐标系与参数}

\subsection{坐标系}

采用圆柱坐标 $(r,\theta,z)$，$z$ 为束轴，右手系，$\theta$ 沿逆时针为正。
评估长度为 $L=1\U{m}$；编织周期长度为 $L_c$，以无量纲周期参数
\begin{equation}
s = z/L_c \in [0,1)
\end{equation}
表征周期内位置。截面指 $z=\text{const}$ 平面，以极坐标 $(r,\theta)$ 描述。

\subsection{基本参数}

\begin{table}[htbp]
\centering
\caption{模型参数与符号}
\label{tab:symbols}
\begin{adjustbox}{max width=\textwidth}
\begin{tabular}{cll}
\toprule
符号 & 含义 & 取值 / 关系 \\
\midrule
$d$ & 细丝直径 & 设计变量（本文取 $0.14\Umm$，与问题二一致） \\
$\rho$ & 细丝半径 & $\rho = d/2$ \\
$t_{\mathrm{ins}}$ & 漆膜等效间隙 & $10\U{\mu m}$ \\
$g$ & 丝间最小中心距（硬约束） & $g = d + t_{\mathrm{ins}} = 0.15\Umm$ \\
$N$ & 丝数 & 设计变量 \\
$M$ & 径向壳层数（不含中心） & $M = \lceil\,\cdot\,\rceil$，由 $N$ 决定 \\
$m$ & 壳层编号 & $m \in \{0,1,\dots,M\}$ \\
$n_m$ & 壳层 $m$ 的丝数 & $\sum_m n_m = N$ \\
$\rho_m$ & 壳层 $m$ 的名义半径 & $\rho_0=0$；$\rho_m = mG$ \\
$G$ & 实际径向间距（环带宽度） & $G = (1+\varepsilon)g$，$\varepsilon\ge 0$ 为膨胀因子 \\
$q$ & 平方因子 & $q = (2M+1)^2$ \\
$A_m$ & 壳层 $m$ 的面积 & $A_0=\pi G^2/4$；$A_m = 2\pi m G^2$ \\
$c_m$ & 壳层 $m$ 的丝数容量 & $c_0=1$；$c_m=\lfloor \pi/\arcsin(g/2mG)\rfloor$ \\
$f_m$ & 壳层 $m$ 的面积占比 & $f_m = A_m / A_{\mathrm{tot}}$ \\
$\ell_{j,m}$ & 丝 $j$ 在壳层 $m$ 的停留弧长 & 由轨迹积分给出 \\
$p_{j,m}$ & 丝 $j$ 的壳层停留占比 & $p_{j,m} = \ell_{j,m}/\sum_{m'}\ell_{j,m'}$ \\
$j$ & 丝编号 & $j \in \{0,\dots,N-1\}$ \\
$k$ & 站点编号 & $k \in \{0,\dots,N-1\}$ \\
\bottomrule
\end{tabular}
\end{adjustbox}
\end{table}

\subsection{统一工况（继承赛题，不得更改）}

无氧铜 $\sigma = 5.8\times10^{7}\U{S/m}$，$\mu_r=1$，$f = 200\U{kHz}$，
$I = 20\U{A}$ 有效值，$J_{dc}\le 4\U{A/mm^2}$，$L=1\U{m}$，
环境 $20\UdegC$，热电解耦，丝间必须电绝缘。
问题四要求与问题二、三在\textbf{同包络、同铜量}下对比，
即包络直径 $\varnothing 3.14\Umm$、铜面积 $5.0953\U{mm^2}$。

\section{编织的数学模型}

\subsection{径向壳层划分}

\begin{definition}[等宽环带划分]
把束截面按半径划分为 $M+1$ 个\textbf{等宽环带} $\{S_m\}_{m=0}^{M}$，宽度均为 $G$：
$S_0$ 为中心圆盘（半径 $G/2$），$S_m$（$m\ge1$）为环带
$\big[(m-\tfrac12)G,\ (m+\tfrac12)G\big]$。壳层面积与面积占比为
\begin{equation}
A_0 = \frac{\pi G^2}{4},\qquad
A_m = \pi\Big[\big(m+\tfrac12\big)^2G^2-\big(m-\tfrac12\big)^2G^2\Big] = 2\pi m G^2 ,
\end{equation}
\begin{equation}
f_m = \frac{A_m}{A_{\mathrm{tot}}},\qquad
A_{\mathrm{tot}} = \sum_{m=0}^{M} A_m = \pi G^2\Big(M(M+1)+\frac14\Big).
\label{eq:Atot}
\end{equation}
\end{definition}

\begin{proposition}[总面积恒等式与面积占比闭式]
\label{prop:Atot}
括号内的量恰为完全平方：
\begin{equation}
\mathcal{D} := M(M+1)+\frac14 = \frac{(2M+1)^2}{4},
\qquad\text{故}\qquad
A_{\mathrm{tot}} = \frac{\pi G^2 (2M+1)^2}{4}.
\end{equation}
记 $q := (2M+1)^2$，则各层面积占比为\textbf{纯有理数且与 $G$ 无关}：
\begin{equation}
\boxed{\;f_0 = \frac{1}{q},\qquad
f_m = \frac{8m}{q}\quad (m\ge1),\qquad \sum_{m=0}^{M}f_m = 1\;}
\label{eq:fm}
\end{equation}
\end{proposition}

\begin{proof}
$4M(M+1)+1 = 4M^2+4M+1 = (2M+1)^2$。由 \eqref{eq:Atot}，
$A_0/(\pi G^2)=\frac{1}{4\mathcal{D}}$，$A_m/(\pi G^2)=2m/\mathcal{D}$；
代入 $\mathcal{D}=q/4$ 得 $f_0=1/q$，$f_m=8m/q$。
求和：$1+\sum_{m=1}^{M}8m = 1+4M(M+1) = (2M+1)^2 = q$。
\end{proof}

\begin{remark}
命题 \ref{prop:Atot} 是全部后续判据的代数基础：\textbf{它把「公平停留」
化为整数分拆问题}——因为 $D_{\mathrm{shell}}$ 只依赖整数丝数 $n_m$ 与
有理数 $Nf_m$，完美性能否精确达到完全由整除性决定。
\end{remark}

\begin{remark}
壳层面积 $\propto m$ 是几何必然：环带面积正比于半径。
这正是「公平停留」要求非均匀停留占比的根源。
\end{remark}

\subsection{单层丝数容量}

等变构造要求每一层容纳整数根丝，且相邻丝心距不小于 $g$。
在半径 $mG$ 的圆上均布 $n$ 个点，相邻弦长为 $2mG\sin(\pi/n)$，
故容量为
\begin{equation}
c_0 = 1,\qquad
c_m = \left\lfloor \frac{\pi}{\arcsin\!\big(g/(2mG)\big)} \right\rfloor
\qquad (m\ge1).
\label{eq:cap}
\end{equation}
$G=g$ 时 $c_m = \lfloor \pi/\arcsin(1/(2m))\rfloor$，即 $6,12,18,25,31,37,\dots$，
在 $m\le3$ 时与六方数 $6m$ 相同、$m\ge4$ 时略大于 $6m$。

\subsection{核心模型：等变编织}

\begin{definition}[编织的等变模型]
称一个编织为\textbf{等变编织}（equivariant braid），若存在
\begin{enumerate}[label=(\alph*), leftmargin=2.4em, itemsep=2pt]
\item 一条\textbf{母轨迹} $\;\gamma:\mathbb{R}\to\mathbb{R}_{\ge0}\times\mathbb{S}^1$，
$\gamma(s) = (r(s),\theta(s))$，满足 $r(s+1)=r(s)$，$\theta(s+1)=\theta(s)+2\pi$；
\item 一组\textbf{平移量} $\delta_j = j/N$，$j=0,\dots,N-1$；
\end{enumerate}
使第 $j$ 根丝的截面轨迹为
\begin{equation}
r_j(z) = r\!\left(\frac{z}{L_c}+\frac{j}{N}\right),\qquad
\theta_j(z) = \theta\!\left(\frac{z}{L_c}+\frac{j}{N}\right)+\frac{2\pi j}{N}.
\label{eq:equivariant}
\end{equation}
\end{definition}

式 \eqref{eq:equivariant} 的含义是：\textbf{每根丝都由同一根母轨迹经「沿轴平移 $jL_c/N$」
与「绕轴旋转 $2\pi j/N$」得到}。这正是「各丝电磁履历严格一致」的数学表达。

\begin{proposition}[等变 $\Rightarrow$ 等履历]
\label{prop:equiv}
在等变模型 \eqref{eq:equivariant} 下，各丝的壳层停留占比相同：
$p_{j,m} = p_{m}$ 与 $j$ 无关。特别地
$D_{\mathrm{shell}}$ 的 $\max_j$ 可去掉，退化为对单一母轨迹的一次积分。
\end{proposition}

\begin{proof}
$r_j(z) = r(z/L_c+j/N)$ 是 $r(\cdot)$ 的平移，而平移不改变函数在任一周期上的
水平集测度：$\{z\in[0,L_c): r_j(z)\in S_m\}$ 的测度与 $j$ 无关。
故 $\ell_{j,m}$ 与 $j$ 无关，从而 $p_{j,m}$ 与 $j$ 无关。
\end{proof}

\subsection{逐站截面形式（便于构造与计算）}

母轨迹的连续形式便于分析，实际构造时用\textbf{逐站截面}离散形式。

\begin{definition}[站点与层序]
把周期离散为 $N$ 个\textbf{站点} $k=0,\dots,N-1$，位于
$z_k = kL_c/N$。定义\textbf{层序}
\begin{equation}
\mathcal{L}:\mathbb{Z}_N\to\{0,1,\dots,M\},\qquad
\#\{k:\mathcal{L}(k)=m\} = n_m\ \ \forall m .
\end{equation}
则第 $j$ 根丝在站点 $k$ 上的状态为
\begin{equation}
\boxed{\;
r_j(z_k) = \rho_{\mathcal{L}((j+k)\bmod N)},\qquad
\theta_j(z_k) = \frac{2\pi\, i_{jk}}{n_{\mathcal{L}((j+k)\bmod N)}} + \psi_{\mathcal{L}((j+k)\bmod N)}
\;}
\label{eq:station}
\end{equation}
其中 $i_{jk}$ 为丝 $j$ 在站点 $k$ 所在壳层内按 $j$ 升序的序号
（$i_{jk}=0,1,\dots,n_m-1$），$\psi_m$ 为壳层整体的方位相位。
\end{definition}

\begin{proposition}[平移生成 $\Rightarrow$ 占据平衡]
\label{prop:occupancy}
在构造 \eqref{eq:station} 下，\textbf{每一个站点}上各壳层的丝数恒为 $n_m$，
与站点 $k$ 无关。
\end{proposition}

\begin{proof}
站点 $k$ 上壳层 $m$ 的丝数为
$\#\{j: \mathcal{L}((j+k)\bmod N)=m\}$。
当 $j$ 遍历 $\{0,\dots,N-1\}$ 时，$(j+k)\bmod N$ 亦遍历 $\mathbb{Z}_N$，
故该计数等于 $\#\{k':\mathcal{L}(k')=m\}=n_m$，与 $k$ 无关。
\end{proof}

命题 \ref{prop:occupancy} 说明：\textbf{平移生成自动保证每一截面都是密排无碰撞的}。
这是「等变」带来的额外好处——公平性与无碰撞性由同一构造同时满足。

\subsection{站点间的 $C^1$ 插值}
\label{sec:interp}

为得到连续可微轨迹，站点之间用三次 Hermite 插值：
对区间 $z\in[z_k,z_{k+1}]$，记 $h=L_c/N$，$t=(z-z_k)/h$，
\begin{equation}
r_j(z) = h_{00}(t)\,r_j(z_k) + h_{10}(t)\,h\,r_j'(z_k)
       + h_{01}(t)\,r_j(z_{k+1}) + h_{11}(t)\,h\,r_j'(z_{k+1}),
\end{equation}
其中 $h_{00}=2t^3-3t^2+1$，$h_{10}=t^3-2t^2+t$，
$h_{01}=-2t^3+3t^2$，$h_{11}=t^3-t^2$，
导数取中心差分 $r_j'(z_k)=\big[r_j(z_{k+1})-r_j(z_{k-1})\big]/(2h)$。
方位角 $\theta_j$ 同法处理，并需先做单调化（消去 $2\pi$ 跳变）。
周期条件 $r_j(z_0)=r_j(z_N)$、$\theta_j(z_N)=\theta_j(z_0)+2\pi$ 保证闭合。

\subsection{生成变换}

模型的全部结构由两个变换生成：

\begin{table}[htbp]
\centering
\caption{生成变换及其性质}
\label{tab:transforms}
\begin{adjustbox}{max width=\textwidth}
\begin{tabular}{clll}
\toprule
变换 & 定义 & 数学性质 & 作用 \\
\midrule
$T_\delta$ & $z \mapsto z + \delta L_c$ & 等距、保测度 & 生成丝间差异（$\delta=j/N$）\\
$R_\alpha$ & $\theta \mapsto \theta + \alpha$ & 等距、保面积 & 生成方位错开（$\alpha=2\pi j/N$）\\
$\Phi_s$ & $\gamma(s) \mapsto \gamma(s+\cdot)$ & 保周期结构 & 编织变换（单参数演化）\\
$\mathcal{E}_m$ & $m \leftrightarrow m\pm1$ & 双射（需留隙） & 层间交换（受定理 \ref{thm:dilation} 约束）\\
\bottomrule
\end{tabular}
\end{adjustbox}
\end{table}

\section{完美编织的公理化定义}

\begin{definition}[完美编织]
\label{def:perfect}
称编织为\textbf{完美编织}，若下列四条同时成立：

\vspace{2pt}
\textbf{(P1) 公平性}（壳层停留与面积成正比）
\begin{equation}
p_{j,m} = f_m = \frac{A_m}{A_{\mathrm{tot}}}
\qquad \forall\, j\in\{0,\dots,N-1\},\ \forall\, m\in\{0,\dots,M\}.
\label{eq:P1}
\end{equation}

\textbf{(P2) 方位各向同性}（方位角覆盖均匀）
\begin{equation}
\int_0^{L_c} e^{\,\mathrm{i}\,q\,\theta_j(z)}\,\mathrm{d}z = 0
\qquad \forall\, q \in \mathbb{Z}\setminus\{0\},\ \forall\, j .
\label{eq:P2}
\end{equation}

\textbf{(P3) 周期闭合}（周期末截面族与周期初全等）
\begin{equation}
\{r_j(L_c),\theta_j(L_c)\}_{j=0}^{N-1}
= \{r_j(0),\theta_j(0)\}_{j=0}^{N-1}
\quad\text{（作为集合相等，允许整体置换）}.
\label{eq:P3}
\end{equation}

\textbf{(P4) 无碰撞}（丝间最小距离不低于硬约束）
\begin{equation}
\min_{j\ne j'}\min_{z\in[0,L_c]}
\big\|\bm{x}_j(z)-\bm{x}_{j'}(z)\big\| \;\ge\; g,
\qquad \bm{x}_j = (r_j\cos\theta_j,\ r_j\sin\theta_j).
\label{eq:P4}
\end{equation}
\end{definition}

\begin{keybox}[为什么 (P1) 能刻画「电磁履历一致」]
邻近效应损耗密度近似正比于 $|H|^2$，而 $|H(r)|\propto r$（束内磁场随半径线性增长）。
故丝 $j$ 的\textbf{平均磁场暴露}为
$\overline{|H|^2}_j \propto \sum_m p_{j,m}\,\overline{r^2}_m$。
若 (P1) 成立，则所有丝的 $\sum_m f_m \overline{r^2}_m$ 相同，
即\textbf{平均磁场暴露严格一致}，从而邻近效应造成的损耗增量一致。
这就是 (P1) 作为「完美」判据的物理依据。
\end{keybox}

\begin{remark}
(P1) 是\textbf{主判据}；(P2) 补足方位方向的公平性（仅 (P1) 不能排除
「始终停在某个方位角」的病态解）；(P3)(P4) 是\textbf{可行性约束}，
不满足则不是可制造的编织。
\end{remark}

\section{评判标准}

\subsection{主判据：换位充分性偏差}

\begin{definition}[换位充分性偏差]
\begin{equation}
D_{\mathrm{shell}}
= \max_{j} \sum_{m=0}^{M}
\left|\, p_{j,m} - f_m \,\right| .
\label{eq:Dshell}
\end{equation}
由命题 \ref{prop:equiv}，等变编织下 $p_{j,m}$ 与 $j$ 无关，故
$D_{\mathrm{shell}} = \sum_m |p_m - f_m|$，一次积分即可算出。
\end{definition}

$D_{\mathrm{shell}}=0$ 当且仅当 (P1) 成立，故它是「完美性」的自然度量。

\subsection{可达到的最优值：闭式与下界}

由命题 \ref{prop:Atot}，$N f_m$ 为有理数；把整数 $n_m$ 尽量贴近 $N f_m$
即得最优。记小数部分 $\phi_m := \{N f_m\} = N f_m - \lfloor N f_m\rfloor$，
$K := \sum_m \phi_m \in \mathbb{Z}$（因 $\sum_m Nf_m = N$ 为整数）。
把「向上取整」的机会分配给 $\phi_m$ 最大的 $K$ 层，即得精确最优值：

\begin{theorem}[一般最优值闭式]
\label{thm:general}
\begin{equation}
\boxed{\;
D_{\mathrm{shell}}^{\min}(M,N)
= \frac{2}{N}\sum_{m\in \mathcal{S}}\big(1-\phi_m\big),
\qquad
\mathcal{S} := \text{$\{\phi_m\}$ 中最小的 $M+1-K$ 个}
\;}
\label{eq:general}
\end{equation}
等价地 $D_{\mathrm{shell}}^{\min} = \dfrac{2}{N}\Big(K-\sum_{m\in\mathcal{T}}\phi_m\Big)$，
$\mathcal{T}$ 为最大的 $K$ 个 $\phi_m$。特别地
\begin{equation}
D_{\mathrm{shell}}^{\min} \;\le\; \frac{2(M+1-K)}{N} \;\le\; \frac{2(M+1)}{N}.
\end{equation}
\end{theorem}

\begin{theorem}[魔数丝数的精确最优值]
\label{thm:closedform}
取六方密排的「魔数」丝数
\begin{equation}
N = 1 + 3M(M+1),\qquad M\in\mathbb{Z}_{\ge1},
\end{equation}
并令 $n_0=1$、$n_m=6m$（$m\ge1$）。则 $D_{\mathrm{shell}}$ 的最优值为
\begin{equation}
\boxed{\;
D_{\mathrm{shell}}^{\min}(N)
= \frac{2M(M+1)}{\big(3M(M+1)+1\big)\big(4M(M+1)+1\big)}
\;}
\label{eq:closedform}
\end{equation}
并满足\textbf{精确恒等式}
\begin{equation}
2N\,D_{\mathrm{shell}}^{\min} \;=\; 1-\frac{1}{(2M+1)^2} \;=\; 1-\frac{1}{q}
\;\xrightarrow[\ M\to\infty\ ]{}\; 1 ,
\qquad\text{且}\qquad
D_{\mathrm{shell}}^{\min} < \frac{1}{2N}.
\label{eq:identity}
\end{equation}
\end{theorem}

\begin{proof}
见附录 \ref{app:proof}。
\end{proof}

\begin{theorem}[精确完美性判据]
\label{thm:exact}
$D_{\mathrm{shell}} = 0$ 当且仅当
\begin{equation}
N = t\,(2M+1)^2 = t\,q,\qquad t\in\mathbb{Z}_{\ge1},
\label{eq:exactcond}
\end{equation}
此时层序 $n_0=t$、$n_m=8mt$ 严格满足 $n_m = N f_m$。
\end{theorem}

\begin{proof}
$D_{\mathrm{shell}}=0$ 要求 $n_m = Nf_m$ 全为整数，即 $q\mid N$（由 \eqref{eq:fm}，
$q f_0=1$、$q f_m=8m$，故 $Nf_m$ 整 $\iff q\mid N$）。
反之 \eqref{eq:exactcond} 下 $Nf_0 = t$、$Nf_m = 8mt$ 均为整数。
\end{proof}

\begin{remark}[完美性与几何可行性是两个独立条件]
定理 \ref{thm:exact} 给出的是\textbf{组合条件}。实际可达还需
$n_m \le c_m$。由于 $c_0=1$，必须有 $t=1$，即 $N=q$；
再由 $m=1$ 层 $n_1=8$ 与 \eqref{eq:cap} 得
\begin{equation}
8 \le c_1 = \left\lfloor\frac{\pi}{\arcsin\big(g/(2G)\big)}\right\rfloor
\quad\Longleftrightarrow\quad
\frac{G}{g} \;\ge\; \frac{1}{2\sin(\pi/8)} = 1.306563\ldots
\quad\Longleftrightarrow\quad
\varepsilon \;\ge\; 0.306563\ldots
\label{eq:eps_exact}
\end{equation}
故\textbf{在 $G=g$（零膨胀）下，$D_{\mathrm{shell}}=0$ 不可能达到}；
但在 $\varepsilon\ge 30.66\%$ 的适度膨胀下可以\textbf{严格达到}。
\end{remark}

\subsection{辅助判据}

\begin{table}[htbp]
\centering
\caption{辅助判据及其判定阈值}
\label{tab:aux}
\begin{adjustbox}{max width=\textwidth}
\begin{tabular}{clll}
\toprule
判据 & 定义 & 理想值 & 说明 \\
\midrule
$D_{\mathrm{az}}$ & $\displaystyle\max_j\Big|\frac{1}{L_c}\int_0^{L_c}e^{\mathrm{i}\theta_j}\mathrm{d}z\Big|$ & $0$ & 方位覆盖均匀度 \\
$D_{\mathrm{per}}$ & 周期末/初截面族的 Hausdorff 距离 & $0$ & 周期闭合性 (P3) \\
$D_{\min}$ & 轨迹上丝间最小距离 & $\ge g$ & 无碰撞 (P4) \\
$C_{\mathrm{env}}$ & $R_{\max}^{\text{编织}} / R_{\max}^{\text{密排}}$ & $\to 1$ & 包络代价 \\
$\kappa_{\min}$ & 轨迹最小曲率半径 & $>0$ & 可制造性 \\
\bottomrule
\end{tabular}
\end{adjustbox}
\end{table}

\subsection{完美性分级}

综合上述判据，给出可直接套用的分级标准：

\begin{table}[htbp]
\centering
\caption{完美性分级判定标准（$q=(2M+1)^2$）}
\label{tab:grade}
\begin{adjustbox}{max width=\textwidth}
\begin{tabular}{clll}
\toprule
等级 & $D_{\mathrm{shell}}$ & 其他要求 & 含义 \\
\midrule
\textbf{理想} & $=0$ & $\varepsilon\ge0.3066$ 且 (P2)(P3)(P4) 全满足 & 严格完美（须适度膨胀）\\
\textbf{近完美} & $\le 1/(2N)$ & (P2)(P3)(P4) 全满足 & 达到魔数族的量化极限 \\
\textbf{合格} & $\le 1/N$ & (P3)(P4) 满足 & 同量级，工程可接受 \\
\textbf{不足} & $> 1/N$ & --- & 需改进层序设计 \\
\textbf{无效} & --- & (P3) 或 (P4) 不满足 & 非可制造编织 \\
\bottomrule
\end{tabular}
\end{adjustbox}
\end{table}

\section{可实现性定理：径向膨胀的必要条件}

判据 (P1) 要求丝在径向层间移动（否则 $p_{j,m}$ 为二值分布，$D_{\mathrm{shell}}$
取到上界）。但层间移动受几何硬约束。以下定理给出\textbf{精确}的必要条件，
其常数由「同一环带上相邻两丝的弦长不得小于 $g$」确定。

\begin{theorem}[径向膨胀定理]
\label{thm:dilation}
设环带宽度为 $G$，丝间最小中心距为 $g$，膨胀因子 $\varepsilon=(G-g)/g$。
则\textbf{存在}满足 (P1)（即 $D_{\mathrm{shell}}=0$）的无碰撞编织，当且仅当
\begin{equation}
\boxed{\;\frac{G}{g} \;\ge\; \frac{1}{2\sin(\pi/8)} = 1.306563\ldots
\qquad\Longleftrightarrow\qquad
\varepsilon \;\ge\; \frac{1}{2\sin(\pi/8)}-1 = 0.306563\ldots\;}
\label{eq:dilation}
\end{equation}
\end{theorem}

\begin{proof}
必要性：由定理 \ref{thm:exact}，$D_{\mathrm{shell}}=0$ 需 $n_m = Nf_m$ 为整数；
又 $c_0=1$ 迫使 $t=1$，$N=q$，于是 $n_1 = 8$。
半径 $G$ 的圆上均布 8 个点，相邻弦长为 $2G\sin(\pi/8)$，
不得小于 $g$，故 $2G\sin(\pi/8)\ge g$，即 $G/g \ge 1/(2\sin(\pi/8))$。

充分性：取 $N=q$、$n_m=8m$（$m\ge1$）、$n_0=1$。
$m=1$ 层是约束最紧的一层（$n_m/(mG) = 8/g$ 与 $m$ 无关，而
$\arcsin$ 单调增使 $\lfloor\pi/\arcsin(\cdot)\rfloor$ 关于 $m$ 不减），
故 $G/g \ge 1/(2\sin(\pi/8))$ 时全部 $n_m\le c_m$，构造可行且 $D_{\mathrm{shell}}=0$。
\end{proof}

\begin{remark}[与六方密排的比较]
六方密排（$G=g$）的层容量 $c_m = 6,12,18,25,\dots$ 与六方数 $6m$ 在 $m\le3$
相同，故 $G=g$ 时最紧的 $m=1$ 层恰为 $6$，只能容纳 $n_1=6$，
于是 $Nf_1 = 8 > 6$，\textbf{公平性 (P1) 必然被违反}——这正是
$D_{\mathrm{shell}}>0$ 的几何根源。
定理 \ref{thm:dilation} 说明：只要把径向间距放大约 $30.66\%$，
即可让最紧的 $m=1$ 层容纳 8 根丝，从而\textbf{严格实现} (P1)。
\end{remark}

\begin{remark}[工程取舍]
定理 \ref{thm:dilation} 给出一个\textbf{定量}结论：
在「同包络」约束下（$G=g$，$\varepsilon=0$），
$D_{\mathrm{shell}}=0$ \textbf{不可达}，只能取到魔数族的
$D_{\mathrm{shell}}^{\min}\approx 1/(2N)$；
而接受 $\varepsilon=30.66\%$ 的包络膨胀（半径比 $1.3066$、截面积比 $1.7071$）
则可\textbf{严格达到} $D_{\mathrm{shell}}=0$。
这一取舍应作为问题四结论的核心论证。
\end{remark}

\section{模型实例}

\subsection{参数表}

\begin{table}[htbp]
\centering
\caption{三种丝数下的模型参数（$d=0.14\Umm$，$g=0.15\Umm$，$G=g$）}
\label{tab:instances}
\begin{adjustbox}{max width=\textwidth}
\begin{tabular}{lccc}
\toprule
项目 & $N=64$（降阶模型） & $N=61$（魔数，$M=4$） & $N=331$（物理模型，$M=10$）\\
\midrule
壳层数 $M$ & 6 & 4 & 10 \\
$q=(2M+1)^2$ & 169 & 81 & 441 \\
最优层序 $\{n_m\}$ & $1,3,6,9,12,15,18$ & $1,6,12,18,24$ & $1,6,12,\dots,60$ \\
是否魔数（$N=1+3M(M+1)$）& 否 & 是 & 是 \\
$D_{\mathrm{shell}}^{\min}$ & $1.942\times10^{-2}$ & $8.096\times10^{-3}$ & $1.507\times10^{-3}$ \\
$1/(2N)$ & $7.813\times10^{-3}$ & $8.197\times10^{-3}$ & $1.511\times10^{-3}$ \\
$2N D_{\mathrm{shell}}^{\min}$ & $2.485$ & $0.9877$ & $0.9977$ \\
分级 & 不足 & 近完美 & 近完美 \\
$\beta = d/\delta$ & $0.947$ & $0.947$ & $0.947$ \\
$\Rac/\Rdc$（仿真） & --- & --- & $4.193$（规则绞合参照）\\
\bottomrule
\end{tabular}
\end{adjustbox}
\end{table}

\begin{remark}
$N=61$ 与 $N=331$ 同为魔数，其 $D_{\mathrm{shell}}$ 精确达到定理
\ref{thm:closedform} 的闭式值，且满足恒等式 $2ND_{\mathrm{shell}}=1-1/q$
（$0.98765$ 与 $0.99773$）；$N=64$ 非魔数，壳层容量无法与面积占比匹配，
$D_{\mathrm{shell}}^{\min}$ 比同量级的 $N=61$ 大 $2.4$ 倍。
\textbf{若丝数可调，应优先选魔数 $N=1+3M(M+1)$}。
\end{remark}

\begin{remark}[若要严格完美]
取 $N=q=(2M+1)^2$ 且 $\varepsilon\ge0.3066$ 即可令 $D_{\mathrm{shell}}=0$。
例如 $M=10$：$N=441$（比 $331$ 多 $33\%$），包络半径 $10.5G = 13.72\,g$
（比六方 $10.5g$ 大 $30.66\%$）。若坚持包络 $\varnothing3.14\Umm$ 不变，
则丝径须按 $1/1.3066$ 缩小至 $d' = 0.1072\Umm$。
\end{remark}

\subsection{模型使用流程}

\begin{enumerate}[label=\textbf{步骤 \arabic*.}, leftmargin=3.6em, itemsep=3pt]
\item \textbf{定参数}：由 $d$、$t_{\mathrm{ins}}$ 得 $g$；给定 $N$，
选取使容量可行（$\sum_m c_m\ge N$）且 $D_{\mathrm{shell}}^{\min}$ 最小的 $M$；
若 $N$ 可自由选择，优先取魔数 $N=1+3M(M+1)$。
\item \textbf{定层序}：由定理 \ref{thm:general} 的闭式（或整数 DP）求最优
$\{n_m\}$，并构造 $\mathcal{L}$ 使层 $m$ 恰占 $n_m$ 个站点（基线为块状排列）。
\item \textbf{生成轨迹}：按式 \eqref{eq:station} 得各站点状态，
按第 \ref{sec:interp} 节插值得 $C^1$ 轨迹。
\item \textbf{算判据}：由轨迹积分得 $p_{j,m}$，按式 \eqref{eq:Dshell} 算
$D_{\mathrm{shell}}$；另算 $D_{\mathrm{az}},D_{\mathrm{per}},D_{\min},C_{\mathrm{env}}$。
\item \textbf{对表分级}：查表 \ref{tab:grade} 得完美性等级；
若未达「近完美」，回到步骤 2 调整层序。
\end{enumerate}

\section{命题的数值验证}
\label{sec:verify}

全部命题已用程序独立验证，采用精确有理算术（\texttt{fractions.Fraction}）与
精确整数动态规划（\texttt{scripts/q4\_def\_verify.py}），无浮点误差。

\begin{table}[htbp]
\centering
\caption{命题验证结果汇总}
\label{tab:verify}
\begin{adjustbox}{max width=\textwidth}
\begin{tabular}{clll}
\toprule
命题 & 检验方式 & 结果 & 结论 \\
\midrule
\ref{prop:Atot} 面积占比闭式
 & $\sum_m f_m = 1$ 对 $M=1\dots7$ & 恒为 $1$（精确有理数）& 精确成立 \\
\ref{prop:equiv} 等变 $\Rightarrow$ 等履历
 & 6 组平移量下比对壳层占比 & 最大偏差 $6.0\times10^{-6}$ & 数值成立 \\
\ref{prop:occupancy} 占据平衡
 & 逐站点统计层占据 & 全部站点恒为 $n_m$，失衡 0 处 & 成立 \\
\ref{thm:general} 一般最优闭式
 & 闭式 vs 整数 DP，$M=1\dots6$，$N$ 全扫 & 全部完全相等（196 组）& 精确成立 \\
\ref{thm:closedform} 魔数闭式
 & $M=1\dots39$ 闭式 vs DP & 全部完全相等 & 精确成立 \\
\ref{thm:closedform} 恒等式
 & $2ND_{\mathrm{shell}} = 1-1/q$，$M=1\dots39$ & 精确成立 & 精确成立 \\
\ref{thm:exact} 精确完美判据
 & 扫描 $N\le700$ 找 $D=0$ & 恰为 $q\mid N$，无一例外 & 精确成立 \\
\ref{thm:dilation} 膨胀阈值
 & 容量可行性二分扫描 & $G/g<1.30656$ 不可行，$\ge$ 可行 & 精确成立 \\
\ref{thm:dilation} 严格完美构造
 & $N=q$ 时 DP 求最优 & $D_{\mathrm{shell}}=0$（$M=1,2,4,5,10$）& 成立 \\
\ref{thm:general} 通用上界
 & $D\le 2(M+1)/N$，$M\le15$ 全扫 & 违反 0 次 & 成立 \\
(P4) 无碰撞
 & 站点内最近邻距离 & 恒为 $0.1500\Umm = g$ & 满足 \\
(P3) 周期闭合
 & $r_j(0)$ vs $r_j(L_c)$ & 最大差 $0$ & 严格闭合 \\
\bottomrule
\end{tabular}
\end{adjustbox}
\end{table}

\section{结论}

\begin{enumerate}[label=(\arabic*), leftmargin=2.4em, itemsep=3pt]
\item \textbf{定义}：完美编织由 (P1)--(P4) 四条公理刻画，其中 (P1) 公平性是
物理核心，(P2) 补方位公平性，(P3)(P4) 为可行性约束。

\item \textbf{模型}：等变编织模型（式 \eqref{eq:equivariant}／\eqref{eq:station}）
由「一条母轨迹 + 平移与旋转两种生成变换」给出，参数明确
（$N,M,\{n_m\},\mathcal{L},\{G\}$），且\textbf{自动满足}等履历（命题 \ref{prop:equiv}）
与占据平衡（命题 \ref{prop:occupancy}）。

\item \textbf{判据}：主判据 $D_{\mathrm{shell}}$（式 \eqref{eq:Dshell}）可计算、可比较；
其最优值有一般闭式 \eqref{eq:general} 与魔数族闭式 \eqref{eq:closedform}，
并满足精确恒等式 $2ND_{\mathrm{shell}}=1-1/q$；通用上界 $D_{\mathrm{shell}}\le 2(M+1)/N$。
分级标准见表 \ref{tab:grade}。

\item \textbf{精确完美}：$D_{\mathrm{shell}}=0$ 当且仅当 $N=t(2M+1)^2$（定理 \ref{thm:exact}）；
结合容量约束 $c_0=1$ 得 $N=(2M+1)^2$，且须 $\varepsilon\ge 30.66\%$
（定理 \ref{thm:dilation}）。故\textbf{完美并非不可达，而是需要付出
约 $30.66\%$ 的径向膨胀代价}；若拒绝该代价（$G=g$），
则最优只能取到 $D_{\mathrm{shell}}\approx 1/(2N)$，即魔数族的近完美。
工程目标因此是\textbf{在包络代价与换位公平性之间定量取舍}。
\end{enumerate}

\appendix

\section{定理 \ref{thm:closedform} 的推导}
\label{app:proof}

取魔数 $N = 1+3A$，$A := M(M+1)$，$n_0=1$，$n_m=6m$，$q=(2M+1)^2$。

\vspace{4pt}
\textbf{第一步：面积占比。} 由 \eqref{eq:fm}，$f_0 = 1/q$，$f_m = 8m/q$。

\vspace{4pt}
\textbf{第二步：目标值与整数丝数之差。} 停留占比 $p_m = n_m/N$：

\emph{中心项}
\begin{equation}
p_0 - f_0 = \frac{1}{N}-\frac{1}{q}
= \frac{q-N}{Nq}.
\end{equation}
由 $q = 4A+1$、$N = 3A+1$，得 $q-N = A$，故
\begin{equation}
\left|p_0-f_0\right| = \frac{A}{Nq}.
\end{equation}

\emph{环带项}（$m\ge1$）
\begin{equation}
p_m - f_m = \frac{6m}{N}-\frac{8m}{q}
= m\cdot\frac{6q-8N}{Nq}.
\end{equation}
由 $6q-8N = 6(4A+1)-8(3A+1) = 24A+6-24A-8 = -2$，故
\begin{equation}
\left|p_m-f_m\right| = \frac{2m}{Nq}.
\end{equation}

\vspace{4pt}
\textbf{第三步：求和。}
\begin{equation}
\sum_{m=1}^{M}\frac{2m}{Nq}
= \frac{2}{Nq}\cdot\frac{M(M+1)}{2}
= \frac{A}{Nq}.
\end{equation}

\vspace{4pt}
\textbf{第四步：合计。}
\begin{equation}
D_{\mathrm{shell}}^{\min}
= \frac{A}{Nq} + \frac{A}{Nq}
= \frac{2A}{Nq}
= \frac{2M(M+1)}{\big(3M(M+1)+1\big)\big(4M(M+1)+1\big)} .
\end{equation}

\vspace{4pt}
\textbf{第五步：恒等式与渐近。} 由上式
\begin{equation}
2N\,D_{\mathrm{shell}}^{\min} = \frac{4A}{q} = \frac{q-1}{q} = 1-\frac{1}{(2M+1)^2},
\end{equation}
其中用到 $q = 4A+1$。因 $q\ge9>1$，故 $2N D_{\mathrm{shell}}^{\min}<1$，
即 $D_{\mathrm{shell}}^{\min} < \dfrac{1}{2N}$（自下方趋近）。
当 $M\to\infty$ 时 $A\sim M^2$、$N\sim3M^2$、$q\sim4M^2$，故
$D_{\mathrm{shell}}^{\min}\sim\dfrac{2M^2}{3M^2\cdot4M^2}=\dfrac{1}{6M^2}\sim\dfrac{1}{2N}$。
证毕。

\vspace{6pt}
\hrule
\vspace{6pt}

\textbf{附录 B：定理 \ref{thm:general} 的说明。}
把 $Nf_m = \lfloor Nf_m\rfloor + \phi_m$ 代入，
最优整数解须把 $\sum \lfloor Nf_m\rfloor$ 与 $N$ 的差 $K=\sum\phi_m$（整数）
通过「向上取整」分配出去；每向上取整一层 $m$，该项偏差为 $1-\phi_m$，
未取整的层偏差为 $\phi_m$。为使总和最小，应选 $\phi_m$ 最大的 $K$ 层向上取整，
得 $\sum_m|n_m-Nf_m| = K - \sum_{\text{top }K}\phi_m$，除以 $N$ 即 \eqref{eq:general}。
数值验证见第 \ref{sec:verify} 节。

\end{document}
